\documentclass[11pt]{article}
\usepackage[utf8]{inputenc}
\usepackage{amsmath,amssymb,amsthm}
\usepackage[margin=1in]{geometry}
\usepackage{hyperref}
\usepackage{xcolor}
\usepackage{enumitem}
\usepackage{newunicodechar}
\newunicodechar{ℝ}{\ensuremath{\mathbb{R}}}
\newunicodechar{⟨}{\ensuremath{\langle}}
\newunicodechar{⟩}{\ensuremath{\rangle}}
\newunicodechar{Γ}{\ensuremath{\Gamma}}
\newunicodechar{·}{\ensuremath{\cdot}}
\newunicodechar{²}{\ensuremath{^2}}
\newunicodechar{ψ}{\ensuremath{\psi}}
\newunicodechar{φ}{\ensuremath{\varphi}}
\newunicodechar{λ}{\ensuremath{\lambda}}
\newunicodechar{≠}{\ensuremath{\ne}}

\setlength{\emergencystretch}{3em}

\newcommand{\userbox}[1]{\par\medskip\begin{quote}\small\textbf{\textcolor{blue!60}{DM:}} #1\end{quote}\medskip}
\newcommand{\claudebox}[1]{\par\medskip\begin{quote}\small\textbf{\textcolor{orange!60!black}{Claude:}} #1\end{quote}\medskip}
\newcommand{\gptbox}[1]{\par\medskip\begin{quote}\small\textbf{\textcolor{green!50!black}{GPT-5.5 Pro:}} #1\end{quote}\medskip}

\title{Tide retrospective:\\\emph{2D cleanup --- refactor through AddSeparable}}
\author{Daniel Murfet (with Claude Opus 4.7)}
\date{6 May 2026}

\begin{document}
\maketitle

\begin{abstract}
A cleanup Tide chained on top of Tide 12 (separable-potential
abstraction). The two existing 2D files
(\texttt{SemiDegenerate.lean}, \texttt{PureQuartic.lean}) are refactored
to apply the new \texttt{AddSeparable} lemmas via small
\texttt{[simp]}-tagged bridge lemmas, instead of re-deriving
\texttt{integral\_prod\_mul} plumbing case-by-case. Net delta is
$-22$ lines across three files (the abstraction itself shed an unused
hypothesis package). Public theorem statements unchanged on the
headline \texttt{cov\_affine\_*} and moment lemmas; only internal proof
structure changes. Build green on first compile, zero
\texttt{sorry}/\texttt{axiom}/\texttt{native\_decide}. The interesting
process moves were two: a deliberate API regression on
\texttt{AddSeparable} (drop hypotheses that Tide 12 added as
discipline; the concrete refactor demonstrated they blocked rather
than helped downstream callers), and a worktree-isolated tide run that
confirmed the structural fix to the recurring concurrent-agent branch
race.
\end{abstract}

\section{Setting}

This is a cleanup Tide rather than a step outward. Tide 12 landed the
\texttt{Laplace.TwoD.AddSeparable} module, which exposes for any
$L(x, y) = U(x) + V(y)$ the generic factorisations
$Z_{2D}(t) = Z_U(t) \cdot Z_V(t)$ and $\langle f \otimes g \rangle =
\langle f \rangle_U \cdot \langle g \rangle_V$. Two existing 2D files
(\texttt{SemiDegenerate.lean}, $L = (\lambda/2) y^2 + x^4/24$, 514
lines; \texttt{PureQuartic.lean}, $L = x^4/24 + y^4/24$, 478 lines)
re-derived these as case-by-case proofs. The cleanup Tide refactors
both files to apply the abstraction's lemmas through a small
\texttt{[simp]} bridge lemma, leaving the public theorem statements
intact.

The Tide chains on top of \texttt{tide/separable-potential} per the
linear-chain branch-hygiene rule; the new branch is
\texttt{tide/twod-cleanup}, off Tide 12's tip
\texttt{725c76b}. The Tide ran in a dedicated \texttt{git worktree} at
\texttt{lean/laplace-cleanup/} --- the first concrete application of
the worktree-isolation fix that three prior retrospectives have
recommended in response to the recurring concurrent-agent branch race.

\section{The thing formalised}

Three pieces, in order:

\paragraph{1. \texttt{AddSeparable.lean} API tightening.}
The \texttt{partitionFunction\_addSeparable\_factor} and
\texttt{integral\_separable\_addSeparable} lemmas drop their
integrability hypotheses, becoming unconditional (matching Mathlib's
\texttt{integral\_prod\_mul}, which is itself unconditional). The
\texttt{gibbsExpectation\_separable\_addSeparable} and
\texttt{gibbsCov\_addSeparable\_fst\_snd\_eq\_zero} lemmas drop the
weight integrability hypotheses but keep the partition-function
nonzero hypotheses (those are load-bearing --- needed to divide).

\paragraph{2. \texttt{SemiDegenerate.lean} refactor.}
A new \texttt{[simp]}-tagged bridge lemma:
\begin{verbatim}
lemma semiDegeneratePotential_eq_addSeparable (lam : ℝ) :
    semiDegeneratePotential lam =
      addSeparable Laplace.OneD.quarticPotential
        (Laplace.OneD.harmonicPotential lam)
\end{verbatim}
Then \texttt{partitionFunction\_factor} and
\texttt{integral\_pow\_pow\_factor} become two-line proofs each:
\begin{verbatim}
theorem partitionFunction_factor (lam t : ℝ) : ... := by
  rw [semiDegeneratePotential_eq_addSeparable]
  exact partitionFunction_addSeparable_factor _ _ _

theorem integral_pow_pow_factor (lam t : ℝ) (m n : ℕ) : ... := by
  simp only [semiDegeneratePotential_eq_addSeparable]
  exact integral_separable_addSeparable _ _ _ (fun x => x ^ m) (fun y => y ^ n)
\end{verbatim}
The \texttt{exp\_neg\_t\_semiDegenerate\_eq\_mul} lemma becomes a
similar two-line corollary of
\texttt{exp\_neg\_t\_addSeparable\_eq\_mul}.

\paragraph{3. \texttt{PureQuartic.lean} refactor.}
Mirror of step 2 with \texttt{pureQuarticPotential\_eq\_addSeparable}
as the bridge.

\paragraph{Public statements unchanged.} The \texttt{cov\_affine\_*}
theorems and the moment lemmas
(\texttt{gibbsExpectation\_fst}/\texttt{snd}/\texttt{fst\_mul\_snd}/\texttt{fst\_sq}/\texttt{snd\_sq})
are unchanged in statement and continue to typecheck through the
refactored proofs. Net delta: $-22$ lines across the three files (64
insertions, 86 deletions).

\section{Proof strategy}

The cleanup pivots on two simple observations:

\begin{enumerate}
\item \textbf{The bridge lemma is \texttt{[simp]}-tagged.}
  The semi-degenerate potential and the corresponding
  \texttt{addSeparable} are pointwise equal as functions $\mathbb{R}
  \times \mathbb{R} \to \mathbb{R}$. Tagging the equation
  \texttt{[simp]} lets a single \texttt{simp only} rewrite the
  integrand of \texttt{integral\_pow\_pow\_factor} from the named
  potential to the \texttt{addSeparable} form, after which the
  AddSeparable lemma applies directly.
\item \textbf{Beta-reduction handles the lambda shapes.}
  \texttt{integral\_separable\_addSeparable U V t f g} expects $f, g :
  \mathbb{R} \to \mathbb{R}$. To recover the existing
  \texttt{integral\_pow\_pow\_factor} statement, we apply at $f = (x
  \mapsto x^m)$, $g = (y \mapsto y^n)$. Lean beta-reduces \texttt{(fun
  x => x ^ m) z.1} to \texttt{z.1 ^ m} automatically.
\end{enumerate}

The API regression on \texttt{AddSeparable.lean} (dropping the
integrability hypotheses) is a separate move: the underlying proof
bodies were already unconditional via Mathlib's
\texttt{integral\_prod\_mul}; the hypotheses had been added in Tide 12
as documentation discipline, not load-bearing. The cleanup Tide's
concrete downstream demonstrates that the hypothesis-bearing form
\emph{blocks} the bridge-rewrite, since the existing
\texttt{partitionFunction\_factor} statements have no positivity
hypotheses to source the integrability witnesses from. The simplest
fix is to drop the discipline.

\section{Roadblocks and resolutions}

\paragraph{The hypothesis package on \texttt{AddSeparable} blocks the
refactor.}
Tide 12 added \texttt{Integrable} hypotheses to
\texttt{partitionFunction\_addSeparable\_factor} and friends, on
GPT-5.5 Pro's recommendation:

\gptbox{For Bochner $\int$, unconditional factorisation statements are
brittle because the theorem you will actually use is typically the
integrable version. \ldots~Downstream files will naturally be able to
apply [the hypothesis-bearing form], since for concrete potentials
they already prove finiteness/integrability of the partition function.}

The cleanup Tide is exactly that downstream, and the prediction did
not hold. The existing \texttt{partitionFunction\_factor} in both 2D
files is \emph{unconditional} (no positivity hypotheses, just $\lambda$
and $t$ as plain reals), so the wrapper has no integrability witnesses
to supply. The choice was either to strengthen the wrapper statements
with positivity (a public API change in two files), or to drop the
hypotheses on the abstraction's lemmas (a public API regression in
one file). Both Claude and GPT (in the deliberation for this Tide)
agreed that dropping the abstraction's hypotheses is the smaller
change; the discipline argument applies in principle, but in practice
the unconditional form is what callers want.

This is worth recording as a Tide-skill calibration: ``Bochner
discipline'' arguments are theoretically right but pragmatically lose
to ``match the existing caller's signature shape'', at least until
Mathlib itself decides to tighten its unconditional integral-prod
lemmas.

\paragraph{Worktree isolation: structural fix lands and works.}
This Tide is the first to run in a dedicated \texttt{git worktree}.
\texttt{lean/laplace-cleanup/} was created from
\texttt{tide/separable-potential}'s tip and held its own
\texttt{HEAD}/index/working-tree throughout the run. The parallel
sessions in \texttt{lean/laplace/} (which had branched onto
\texttt{tide/bounded-prior-continuous-test} and back to \texttt{main}
several times during this session) did not interfere. Three previous
retrospectives recommended this fix; this is the first concrete
confirmation. Worth lifting to a Tide-skill preflight step.

\section{What was Mathlib, what was new}

\paragraph{Mathlib provided.} Nothing new. The cleanup adds no new
Mathlib API consumption.

\paragraph{The seabed (Tide 12) provided.}
The five generic factorisation lemmas of the
\texttt{AddSeparable} module --- the definition, the Boltzmann
decomposition, partition factorisation, integral factorisation,
expectation factorisation, and the mixed-covariance vanishing.

\paragraph{What was new in this Tide.}
The two \texttt{[simp]} bridge lemmas\footnote{Names:
\texttt{semiDegeneratePotential\_eq\_addSeparable} and
\texttt{pureQuarticPotential\_eq\_addSeparable}.} identifying the
existing 2D potentials as instances of \texttt{addSeparable}.
Everything else is either deletion or a syntactic re-targeting of
existing proof bodies through the bridge.

\section{Lessons}

\begin{itemize}
\item \emph{Worktree isolation is the right structural fix.} Three
prior retrospectives recommended this; one Tide actually used it; it
worked. Promote to a Tide-skill preflight (``run the Tide in a
\texttt{git worktree} when concurrent sessions exist on the seabed
repo'').

\item \emph{The Bochner-discipline argument loses on first contact
with downstream.} Tide 12 added \texttt{Integrable} hypotheses to
factorisation lemmas as documentation discipline. Tide 13 (this one)
removed them, because they blocked the cleanest refactor of two
unconditional callers. The right convention for Mathlib-shaped
factorisation lemmas with unconditional underlying API
(\texttt{integral\_prod\_mul}) is: stay unconditional unless a
specific caller's needs strengthen the case. ``Discipline'' alone is
not enough; concrete usage is.

\item \emph{Cleanup Tides have higher leverage than they look.}
$-22$ lines is small. But the cleanup makes A2 (quartic--sextic mixed)
and A4 (mixed even-power 2D moments) one-line specialisations rather
than file-sized re-derivations. The next Tide that adds a 2D
potential pays a constant cost, not a file cost.
\end{itemize}

\section{Follow-ups}

\begin{itemize}
\item \emph{Quartic--sextic mixed potential (A2 from the 6 May
survey).} Now a clean specialisation: define
\texttt{quarticSexticPotential = addSeparable quarticPotential
sexticPotential}, supply a one-line bridge, and the partition + moment
factorisations come for free.

\item \emph{Mixed even-power 2D moments (A4).} One-line corollary of
the integral-factorisation lemma at $f = x \mapsto x^m,\, g = y
\mapsto y^n$, in either of the two refactored files (or generically
over \texttt{addSeparable}).

\item \emph{Promote 2D atom integrability to \texttt{AddSeparable.lean}.}
Both files still carry private \texttt{<name>\_integrable\_pow\_pow}
witnesses built via \texttt{Integrable.mul\_prod} from two 1D
witnesses. A generic
\texttt{integrable\_separable\_addSeparable} in
\texttt{AddSeparable.lean} would shrink each downstream file by
another $\sim$30 lines. GPT flagged this as a strong optional add-on;
deferred to a follow-up Tide.

\item \emph{Generic \texttt{gibbsCov} API (orthogonal future Tide).}
Still the natural infrastructure Tide for the affine-covariance
Candidate B from Tide 12's deliberation\footnote{Tentative names:
\texttt{gibbsCov\_symm}, \texttt{gibbsCov\_add\_left},
\texttt{gibbsCov\_const\_left}, \texttt{gibbsCov\_smul\_left} and
right-hand variants.}. Independent of this cleanup but unblocked by it.

\item \emph{Promote worktree isolation to the Tide skill's preflight.}
The skill currently flags the concurrent-agent race as a recurring
roadblock but does not prescribe a fix. Three retrospectives have
named the same fix; one Tide has now shown it works. Worth landing.
\end{itemize}

\section*{Acknowledgements}

GPT-5.5 Pro's deliberation contributed both the recommendation to
keep public theorem statements unchanged (rejecting Candidate B's
definitional refactor) and the architectural decision to drop the
integrability hypotheses on the AddSeparable lemmas rather than
strengthen the wrapper statements with positivity. The full consult is
preserved in the primer project's tide-log directory under the SRI
repo.

\end{document}
